\begin{lemma}
\label{lemma-canonical-class}
Let $\mathcal{C}$ be a site. Let $\mathcal{O} \to \mathcal{O}_0$
be a surjection of sheaves of rings whose kernel is an ideal sheaf
$\mathcal{I}$ of square zero. For every object
$K_0$ in $D^-(\mathcal{O}_0)$ there is a canonical map
$$
\omega(K_0) :
K_0 \longrightarrow
K_0 \otimes_{\mathcal{O}_0}^\mathbf{L} \mathcal{I}[2]
$$
in $D(\mathcal{O}_0)$ such that for any map
$K_0 \to L_0$ in $D^-(\mathcal{O}_0)$ the diagram
$$
\xymatrix{
K_0 \ar[d] \ar[rr]_-{\omega(K_0)} & &
(K_0 \otimes^\mathbf{L}_{\mathcal{O}_0} \mathcal{I})[2] \ar[d] \\
L_0 \ar[rr]^-{\omega(L_0)} & &
(L_0 \otimes^\mathbf{L}_{\mathcal{O}_0} \mathcal{I})[2]
}
$$
commutes.
\end{lemma}

\begin{proof}
Represent $K_0$ by any complex
$\mathcal{K}_0^\bullet$ of $\mathcal{O}_0$-modules.
Apply Lemma \ref{lemma-lift-complex}
with $\mathcal{G}^n = 0$ for all $n$.
Denote $d : \mathcal{F}^n \to \mathcal{F}^{n + 1}$
the maps produced by the lemma. Then we see that
$d \circ d : \mathcal{F}^n \to \mathcal{F}^{n + 2}$
is zero modulo $\mathcal{I}$. Since $\mathcal{F}^n$ is flat,
we see that
$\mathcal{I}\mathcal{F}^n =
\mathcal{F}^n \otimes_{\mathcal{O}} \mathcal{I} =
\mathcal{F}^n_0 \otimes_{\mathcal{O}_0} \mathcal{I}$.
Hence we obtain a canonical map of complexes
$$
d \circ d : \mathcal{F}_0^\bullet
\longrightarrow
(\mathcal{F}_0^\bullet \otimes_{\mathcal{O}_0} \mathcal{I})[2]
$$
Since $\mathcal{F}_0^\bullet$ is a bounded above complex
of flat $\mathcal{O}_0$-modules, it is K-flat and may be used
to compute derived tensor product. Moreover, the map of
complexes $\mathcal{F}_0^\bullet \to \mathcal{K}_0^\bullet$
is a quasi-isomorphism by construction. Therefore the source and
target of the map just constructed represent $K_0$ and
$K_0 \otimes_{\mathcal{O}_0}^\mathbf{L} \mathcal{I}[2]$
and we obtain our map $\omega(K_0)$.

\medskip\noindent
Let us show that this procedure is compatible with maps of complexes.
Namely, let $\mathcal{L}_0^\bullet$ represent another object of
$D^-(\mathcal{O}_0)$ and suppose that
$$
\mathcal{K}_0^\bullet \longrightarrow \mathcal{L}_0^\bullet
$$
is a map of complexes. Apply Lemma \ref{lemma-lift-complex}
for the complex $\mathcal{L}_0^\bullet$, the flat modules
$\mathcal{F}^n$, the maps $\mathcal{F}^n \to \mathcal{F}^{n + 1}$,
and the compositions
$\mathcal{F}^n \to \mathcal{K}_0^n \to \mathcal{L}_0^n$
(we apologize for the reversal of letters used).
We obtain flat modules $\mathcal{G}^n$, maps
$\mathcal{F}^n \to \mathcal{G}^n$, maps
$\mathcal{G}^n \to \mathcal{G}^{n + 1}$, and maps
$\mathcal{G}^n \to \mathcal{L}_0^n$ with all properties
as in the lemma. Then it is clear that
$$
\xymatrix{
\mathcal{F}_0^\bullet \ar[d] \ar[r] &
(\mathcal{F}_0^\bullet \otimes_{\mathcal{O}_0} \mathcal{I})[2] \ar[d] \\
\mathcal{G}_0^\bullet \ar[r] &
(\mathcal{G}_0^\bullet \otimes_{\mathcal{O}_0} \mathcal{I})[2]
}
$$
is a commutative diagram of complexes.

\medskip\noindent
To see that $\omega(K_0)$ is well defined, suppose that we have two complexes
$\mathcal{K}_0^\bullet$ and $(\mathcal{K}'_0)^\bullet$
of $\mathcal{O}_0$-modules representing $K_0$ and two systems
$(\mathcal{F}^n, d : \mathcal{F}^n \to \mathcal{F}^{n + 1},
\mathcal{F}^n \to \mathcal{K}_0^n)$
and
$((\mathcal{F}')^n, d : (\mathcal{F}')^n \to (\mathcal{F}')^{n + 1},
(\mathcal{F}')^n \to \mathcal{K}_0^n)$
as above. Then we can choose a complex $(\mathcal{K}''_0)^\bullet$
and quasi-isomorphisms
$\mathcal{K}_0^\bullet  \to (\mathcal{K}''_0)^\bullet$
and
$(\mathcal{K}'_0)^\bullet  \to (\mathcal{K}''_0)^\bullet$
realizing the fact that both complexes represent $K_0$ in the
derived category. Next, we apply the result of the previous paragraph
to
$$
(\mathcal{K}_0)^\bullet \oplus (\mathcal{K}'_0)^\bullet
\longrightarrow
(\mathcal{K}''_0)^\bullet
$$
This produces a commutative diagram
$$
\xymatrix{
\mathcal{F}_0^\bullet \oplus (\mathcal{F}'_0)^\bullet
\ar[d] \ar[r] &
(\mathcal{F}_0^\bullet \otimes_{\mathcal{O}_0} \mathcal{I})[2] \oplus
((\mathcal{F}'_0)^\bullet \otimes_{\mathcal{O}_0} \mathcal{I})[2] \ar[d] \\
\mathcal{G}_0^\bullet \ar[r] &
(\mathcal{G}_0^\bullet \otimes_{\mathcal{O}_0} \mathcal{I})[2]
}
$$
Since the vertical arrows give quasi-isomorphisms on the summands
we conclude the desired commutativity in $D(\mathcal{O}_0)$.

\medskip\noindent
Having established well-definedness, the statement on compatibility
with maps is a consequence of the result in the second
paragraph.
\end{proof}
